\subsection{Quy tắc nghiệp vụ trọng yếu và Ánh xạ dữ liệu}

\subsubsection{Các quy tắc nghiệp vụ (Business Rules)}
Các quy tắc nghiệp vụ dưới đây có tác động trực tiếp đến việc xác định các thực thể, bản số mối quan hệ và các ràng buộc toàn vẹn trong ERD.

\begin{description}
    \item[BR-1] \textbf{Ràng buộc sở hữu Chiến dịch:} Một Doanh nghiệp (Merchant) có thể quản lý nhiều Chiến dịch (Campaign). Tuy nhiên, mỗi Chiến dịch bắt buộc phải thuộc về duy nhất một Doanh nghiệp (Quan hệ 1:N).
    \item[BR-2] \textbf{Ràng buộc độc lập Số dư:} Số dư (Credit Balance) của một Khách hàng (Customer) được tính độc lập trên từng Merchant. Một Customer có thể có nhiều ví điểm tương ứng với các Merchant khác nhau, nhưng mỗi cặp (CustomerID, MerchantID) có tối đa một ví; cặp này phải UNIQUE.
    \item[BR-3] \textbf{Quy tắc Bất biến Giao dịch (Ledger Immutability):} Bất kỳ sự thay đổi số dư nào đều sinh ra một bản ghi giao dịch (Credit Transaction) mới. Hệ thống không cho phép sửa đổi hoặc xóa giao dịch đã hoàn tất.
    \item[BR-4] \textbf{Ràng buộc Miền giá trị (Domain Constraints):} Số dư khả dụng không âm; PointsMultiplier dương. Amount là biến động có dấu: Earn dương, Redeem âm, Refund khác 0 và đảo dấu biến động gốc. Tầng dịch vụ phải kiểm tra giao dịch gốc khi hoàn điểm. Campaign.Status thuộc Draft/Active/Expired; Merchant.Status thuộc Active/Suspended. EndDate có thể NULL, nếu có phải không trước StartDate. ActionType là mã hành động có thể mở rộng, không giới hạn ở các ví dụ.
    \item[BR-5] \textbf{Ràng buộc Mã băm Blockchain:} Giao dịch đã đồng bộ On-chain lưu mã băm đối soát \texttt{SuiTxDigest}, duy nhất khi khác NULL; nhiều giao dịch chưa đồng bộ được phép NULL. Phạm vi đồ án giả định một digest ứng với tối đa một giao dịch nội bộ, không mô hình hóa SuiObject riêng. Nếu một giao dịch on-chain chứa nhiều sự kiện nghiệp vụ thì cần mở rộng khóa đối soát, không dùng giả định này.
    \item[BR-6] \textbf{Tài khoản và định danh:} Mỗi User có đúng một Role tại một thời điểm; Email đăng nhập không NULL và duy nhất toàn hệ thống. ManagerID tùy chọn, tham chiếu một User khác. ApiKey không NULL và duy nhất. WalletAddress tùy chọn, duy nhất khi đã được cung cấp.
\end{description}

\subsubsection{Ánh xạ Yêu cầu vào Đối tượng Dữ liệu chính}
Dựa trên các yêu cầu chức năng (FR-DB) và nhóm dữ liệu cần quản lý, bảng dưới đây tổng hợp ánh xạ sang các thực thể CSDL cốt lõi:

\vspace{0.3cm}
\begin{table}[H]
    \centering
    \caption{Ánh xạ Yêu cầu và Nhóm dữ liệu vào Thực thể}
    \label{tab:anh_xa_du_lieu}
    \vspace{0.2cm}
    \renewcommand{\arraystretch}{1.3}
    \begin{tabular}{|>{\raggedright\arraybackslash}p{3cm}|>{\raggedright\arraybackslash}p{3cm}|>{\raggedright\arraybackslash}p{8cm}|}
        \hline
        \textbf{Nhóm Dữ liệu} & \textbf{Mã Yêu cầu} & \textbf{Thực thể quản lý \& Vai trò} \\
        \hline
        Tổ chức \& Dân cư & FR-DB-01, 03 & \texttt{Tenant}, \texttt{Merchant}, \texttt{Customer}: Quản lý định danh và hồ sơ đối tác, khách hàng. \\
        \hline
        Phân quyền & FR-DB-02 & \texttt{User}, \texttt{Role}, \texttt{Permission}: Kiểm soát truy cập RBAC. \\
        \hline
        Loyalty & FR-DB-04, 06 & \texttt{Credit}, \texttt{Campaign}, \texttt{RewardRule}: Quy tắc tích điểm, đổi thưởng và thời hạn. \\
        \hline
        Giao dịch On/Off & FR-DB-05, 07 & \texttt{CreditTransaction}: Ghi nhận biến động số dư và tham chiếu mã băm Blockchain. \\
        \hline
        Bảo mật \& Audit & FR-DB-08, 09 & \texttt{APICredential}, \texttt{AuditLog}, \texttt{FraudEvent}: Quản lý API Key, nhật ký hệ thống và chống gian lận. \\
        \hline
        Thương mại & FR-DB-10 & \texttt{PricingPlan}, \texttt{Subscription}, \texttt{Invoice}: Quản lý gói cước và thanh toán của Merchant. \\
        \hline
    \end{tabular}
\end{table}

\textit{Phạm vi bảng ánh xạ:} Permission, AuditLog, FraudEvent, PricingPlan, Subscription và Invoice chỉ thuộc yêu cầu mở rộng; không có trong ERD, lược đồ logic hay DDL của đồ án. Các mã FR-DB ở bảng này dùng để nhóm yêu cầu dữ liệu, không phải mã FS của bảng chức năng.
